% 06-even.tex -- the even family.  Formulas F14-F17.
% Input by mostlyharmless.tex.  No preamble here.

\section{The even family}
\label{sec:even}

For $l+m$ even the parity dichotomy of section~\ref{sec:parity} gives no factor of
$\sqrt{1-\rho^2}$, and the radial integral does not collapse to a single spherical
Bessel function.  It becomes a finite sum of ordinary Bessel functions instead, with
coefficients that are exact rationals.

\subsection{The mode as a polynomial in \texorpdfstring{$\rho^2$}{rho squared}}
\label{sec:even-polynomial}

With $u = \rho^2$ and $x = \sqrt{1-u} = \cos\theta$, the associated Legendre
function is
\begin{equation}
  P_{lm}\!\left(\sqrt{1-u}\right) = u^{m/2} A(u) , \qquad
  A(u) = \left. \frac{d^m P_l}{dx^m} \right|_{x^2 = 1-u} ,
  \label{eq:even-A}
\end{equation}
in the convention of section~\ref{sec:conventions}, which carries no Condon--Shortley
phase.  For $l+m$
even the $m$-th derivative is an even function of $x$, so it is a polynomial in
$x^2 = 1-u$, and $A$ is an ordinary polynomial.  That is the whole difference
between the families: the odd family has a half-integer power to exploit in its
doubling, and here its absence is what forces a sum.

\subsection{The Zernike radial polynomials}
\label{sec:zernike}

These are the natural basis for functions of $\rho$ with a fixed parity.  With the
normalisation $R_n^m(1) = 1$, and $n-m$ even,
\begin{equation}
  R_n^m(\rho) = \sum_{s=0}^{(n-m)/2} (-1)^s
  \frac{(n-s)!}{s! \left(\frac{n-m}{2}+m-s\right)! \left(\frac{n-m}{2}-s\right)!}\,
  \rho^{n-2s} .
  \label{eq:zernike-sum}
\end{equation}
Equivalently, and better for evaluation,
\begin{equation}
  R_n^m(\rho) = (-1)^{(n-m)/2} \rho^m P_{(n-m)/2}^{(m,0)}\!\left(1-2\rho^2\right) ,
  \label{eq:zernike-jacobi}
\end{equation}
a Jacobi polynomial in $1-2\rho^2$.

They are orthogonal against the weight $\rho$,
\begin{equation}
  \int_0^1 R_n^m(\rho)\, R_{n'}^m(\rho)\, \rho\, d\rho
  = \frac{\delta_{nn'}}{2n+2} ,
  \label{eq:zernike-orthogonality}
\end{equation}
\begin{itemize}
  \item[$\checkmark$] \code{check_zernike_jacobi} relates \eqref{eq:zernike-sum} and
    \eqref{eq:zernike-jacobi} by expanding both as polynomials, for $n \le 8$.
    \\[1pt] \vtag{8}\quad \emph{Exact}
  \item[$\checkmark$] \code{check_zernike_orthogonality} checks
    \eqref{eq:zernike-orthogonality} in rational arithmetic for every pairing, for $n \le 8$.
    \\[1pt] \vtag{9}\quad \emph{Exact}
\end{itemize}

\subsection{The expansion, and its coefficients}
\label{sec:even-projection}

Write the mode's radial part as $\rho^m A(\rho^2)$, a polynomial of degree $l$, and
expand it in Zernike polynomials of the same order,
\begin{equation}
  \rho^m A(\rho^2) = \sum_{k=0}^{(l-m)/2} c_k\, R_{m+2k}^m(\rho) .
  \label{eq:even-expansion}
\end{equation}
The expansion is finite because both sides have the same degree and parity, and
\eqref{eq:zernike-orthogonality} turns the coefficients into integrals,
\begin{equation}
  c_k = 2(m+2k+1) \int_0^1 \rho^m A(\rho^2)\, R_{m+2k}^m(\rho)\,
  \rho\, d\rho .
  \label{eq:even-coefficients}
\end{equation}
The integrand is a polynomial, so every $c_k$ is an exact rational.  They are not
tabulated: they are computed on demand and memoised (decision D19), which is why
there is no data file and no maximum $l$.

\begin{itemize}
  \item[$\checkmark$] \code{check_zernike_projection} obtains the same coefficients a second
    way, by matching powers of $\rho$ from the highest downwards.  That system is triangular,
    since the highest power of the mode occurs in $R_{l,m}$ alone, and the two routes agree,
    in rational arithmetic, for $l \le 6$.
    \\[1pt] \vtag{12}\quad \emph{Exact}
\end{itemize}

\subsection{The Zernike--Hankel identity}
\label{sec:zernike-hankel}

The integral of a Zernike polynomial against a Bessel function is elementary
\citep{zernike1934,bornwolf}:
\begin{equation}
  \int_0^1 R_n^m(\rho)\, J_m(a\rho)\, \rho\, d\rho
  = (-1)^{(n-m)/2} \frac{J_{n+1}(a)}{a} .
  \label{eq:zernike-hankel}
\end{equation}
This is the identity that converts the expansion into Bessel functions, so its sign
and normalisation matter.  Certificate \vtag{10} expands both sides in powers of $a$
--- the left from the exact moments of $R_n^m$, the right from the Bessel series ---
and compares every power up to $a^{n+2J}$.  \emph{Exact}, with that stated series
scope, $n \le 6$ and $J = 6$.

\subsection{The result}
\label{sec:even-result}

Substituting \eqref{eq:even-expansion} into the radial integral and applying
\eqref{eq:zernike-hankel} term by term gives
\begin{equation}
  X_{lm}(q) = \sum_{k=0}^{(l-m)/2} c_k\, (-1)^k\,
  \frac{J_{m+2k+1}(2\pi q)}{2\pi q} , \qquad l+m \text{ even} ,
  \label{eq:even-radial}
\end{equation}
a finite sum of ordinary Bessel functions of odd orders $m+1, m+3, \dots, l+1$.  For
$l = m = 0$ it is $J_1(2\pi q)/(2\pi q)$, which is the uniform disk of
section~\ref{sec:uniform-disk}, so that anchor is recovered here by a different
route.

At $q = 0$ only the term with $m+2k = 0$ survives, since $J_\nu(a) \approx a^\nu$
there:
\begin{equation}
  X_{lm}(0) = \begin{cases}
    c_0/2 , & m = 0 , \\[2pt]
    0 , & m \ge 2 ,
  \end{cases}
  \label{eq:even-limit}
\end{equation}
the case $m = 1$ being absent from the even family.

\subsection{The coefficients to \texorpdfstring{$l \le 3$}{l <= 3}}
\label{sec:even-table}

The coefficients are computed in exact rational arithmetic and emitted as LaTeX by
\code{make data}, so the numbers below are derived rather than transcribed.

\IfFileExists{generated/coefficients.tex}{\begin{center}
\begin{tabular}{llll}
\hline
$l$ & $m$ & $k$ & $c_k$ \\
\hline
$0$ & $0$ & $0$ & $1$ \\
$1$ & $1$ & $0$ & $-1$ \\
$2$ & $0$ & $0$ & $\frac{1}{4}$ \\
$2$ & $0$ & $1$ & $- \frac{3}{4}$ \\
$2$ & $2$ & $0$ & $3$ \\
$3$ & $1$ & $0$ & $-1$ \\
$3$ & $1$ & $1$ & $\frac{5}{2}$ \\
$3$ & $3$ & $0$ & $-15$ \\
\hline
\end{tabular}
\end{center}
}{%
  \emph{Not generated yet --- run \code{make data}.}}

\subsection{What was checked}
\label{sec:even-register}

\begin{center}
\begin{tabular}{lll}
\hline
tag & claim & rung \\
\hline
\vtag{8} & \eqref{eq:zernike-sum} against \eqref{eq:zernike-jacobi} & exact, $n \le 8$ \\
\vtag{9} & orthogonality \eqref{eq:zernike-orthogonality} & exact, $n \le 8$ \\
\vtag{10} & the Hankel identity \eqref{eq:zernike-hankel}, by power series & exact, $n \le 6$, powers to $a^{18}$ \\
\vtag{12} & \eqref{eq:even-coefficients} against coefficient matching & exact rational, $l \le 6$ \\
\vtag{14} & \eqref{eq:even-radial} against the quadrature, $q = 0$ included & numerical, $l \le 6$ \\
\vtag{22} & the Zernike transform identity, against the library's $R_n^m$ & numerical, $n \le 8$ \\
\hline
\end{tabular}
\end{center}

Certificate \vtag{22} came with the disk-mode proposal that has since been retired: the
integrand of the Zernike transform is a polynomial times a
Bessel function, so Gauss--Legendre quadrature evaluates it to machine precision, and any
residual is the identity's own rather than the quadrature's.  It runs over every $(n,m)$ with
$n \le 8$ and $n - m$ even, against the library's own $R_n^m$, so it certifies the convention
as well as the formula.
